% 迁移来源：NSFC_2026_CIESM-STSR/extraTex/3.1.研究基础.tex
\subsubsection{与本项目相关的研究工作积累}

申请人\CiesmValue{applicant}为\CiesmValue{institution}\CiesmValue{professional-title}、\CiesmValue{degree}、\CiesmValue{mentor}，已开展多尺度时间序列、能源预测和多源信息融合研究。这些研究积累主要包括从多时间尺度中提取可复用表征、分析不同信息源之间的关系，以及组织模型训练和比较实验，可用于本课题的资料处理、联合表征和时间连续性分析。

申请人在 DWT-Former 工作中担任通讯作者，研究小波多尺度特征与 Transformer 时间表示的结合，成果发表于 Energy\cite{DWT2025}。该工作为本课题分解缓慢变化背景和局地波动提供了技术基础。申请人在 WaveKAN 工作中担任共同第一作者，研究自适应小波分解与长时序建模，成果发表于 Neurocomputing\cite{WaveKAN2026}。该工作为设计不同时间尺度的对照及分析长期变化提供了方法基础。

上述成果与 CIESM 气候场的研究对象不同，可借鉴之处主要在于模型组织与实验方法。为检验本课题的核心设计，申请人已完成一项先行试验，结果见下节。

\subsubsection{已完成的先行试验及其结论}

先行试验在公开 ERA5 资料上复现本课题的两因素对照设计，检验这一设计能否回答“地形条件与多变量共享表征各自贡献多少”。\textbf{该试验为方法学验证，不等同于 CIESM 正式任务}：粗场由 ERA5 0.25\textdegree{} 经面积加权算子 $A$ 自粗化到 1.0\textdegree{}（线性 4 倍、面积 16 倍），属同源配对，不含 CIESM 的域差距与非配对迁移问题。

\paragraph{设置}
区域取 73.00—95.75\textdegree{}E、34.25—50.00\textdegree{}N，细网格 64$\times$92、粗网格 16$\times$23；时段 2018—2020 年逐日 06 UTC 共 1096 天，按连续时段划分：2018—2019 年用于训练（730 天），2020 年上半年用于验证（182 天），下半年用于测试（184 天）。变量为近地面气温、10\,m 风速（由风分量逐时刻合成）、地表净短波与降水，地形通道为高程、次网格地形起伏、坡度与陆海掩膜。粗化算子经校验：细粗场面积加权均值最大差 $9.4\times10^{-7}$\,K，块常数展开经 $A$ 回聚合残差恰为零。独立与联合模型按等预算反解匹配参数量，无地形和有地形配置的差异分别约为 0.13\% 和 2.13\%；每配置 3 个随机种子。图中误差为各随机种子 MAE 的平均值；下述增益及其区间则先对 3 个种子的预测取平均，再计算 MAE 差，区间由 7 天时间块 bootstrap 给出。

\paragraph{结论一：地形条件降低了气温重建误差}
独立模型加入地形后，气温平均绝对误差为 0.441\,K，较双线性插值的 1.641\,K 降低约 73\%。在独立模型内部比较有无地形配置，按种子预测平均后的统计口径，地形条件带来的 MAE 增益为 0.217\,K，95\% 区间为 $[+0.210,+0.225]$，高于事前写定的逐变量最小有意义增益 0.150\,K，据此支持本试验中的气温地形增益。这一结果为地形信息有助于恢复气温亚网格结构提供了先行证据。

\paragraph{结论二：等信息对照未发现共享表征的有效增益}
在无地形配置下，独立与联合模型的输入信息相同、参数量差异约为 0.13\%。按种子预测平均后的统计口径，共享表征相对独立基线的气温 MAE 增大 0.020\,K，其余三变量的增益均低于各自的最小有意义增益，跨变量秩相关矩阵误差由 0.112 升至 0.121。有地形联合模型的种子平均气温 MAE 为 0.475\,K，较插值降低约 71\%，但这一总改善不能归因于共享表征。等信息独立基线使地形条件与联合架构的贡献能够分别检验，避免将二者混在一起解释。
图~\ref{fig:pilot-error} 展示了等信息对照下四变量的重建误差。
\CiesmFigure{figures/F08-pilot-error.pdf}{先行试验：等信息对照下四变量重建误差。柱高为测试期 3 个随机种子的 MAE 均值，误差棒为种子间标准差}{fig:pilot-error}

\paragraph{结论三：未加显式投影时仍有明显的聚合残差}
块常数插值的聚合残差恰为零，验证了检查方法的正确性。学习模型的气温聚合残差为 0.058—0.079\,K，约占其自身误差的 13\%—18\%，双线性插值的聚合残差达到 0.845\,K。降水聚合残差沿天山山脉走向呈系统性负偏，具有明确的空间分布特征。这些实测结果说明，有必要在研究内容三的 TerraSR-XJ 中引入显式投影。
图~\ref{fig:pilot-agg} 展示了重建场回聚合后与粗输入之差。
\CiesmFigure{figures/F09-pilot-aggregation.pdf}{先行试验：重建场经算子 $A$ 回聚合后与粗输入之差，未施加显式投影时的一致性基线}{fig:pilot-agg}

\paragraph{先行试验的局限与阈值使用条件}
本试验尚未涉及 CIESM 域差距与非配对迁移，也未开展目标场信息含量检验和台站独立验证；正式研究设定的空间结构主指标同样尚未检验。这些问题将由研究内容一至四处理；面积 100 倍倍率下结论是否成立，也需正式实验确认。试验还表明，逐变量最小有意义增益的绝对下限必须与时间聚合口径相适应：将日总量口径的降水下限 0.30\,mm 直接用于逐时次口径，会使阈值比误差本身大一个量级。正式研究先将降水统一为日总量，再使用申请书写定的阈值规则；本次逐时次试验的降水判定不外推为日尺度结论。

先行试验的数据处理过程、模型配置、评价指标与图件均已归档，可支持结果复现与核查。

\subsubsection{尚待开展的验证工作}

目前尚无针对 CIESM 的数据产品、训练模型、消融实验或性能记录，因此本申请不报告已取得的 CIESM 重建指标。第一季度将通过数据核查与基线实验，检验上述先行试验的设计能否用于真实模式资料。

\subsubsection{已取得的研究工作成绩}

两项代表性成果体现了申请人在多尺度时间表征和自适应序列建模方面的积累，可用于本课题的时空信息组织、模型比较和训练稳定性分析。完整题录见参考文献，既有成果与本课题的气候实验分别列示。除上述论著外，申请人已围绕本课题完成前述先行试验，形成可复现的数据处理流程、四配置受控对照结果与两项诊断图件，并验证了对照设计能够区分不同因素的贡献。

申请人主持在研的“融合多源数据大模型驱动的光伏发电功率预测方法研究”，并参与“数据驱动下基于互补型生成对抗网络与神经统计模型的公路沙害智能预警研究”，积累了多源数据整理和区域相关资料处理经验。项目概况及与本课题的区别见下文“相关在研项目”。已有项目的任务、资金和成果归属按各自任务书执行，本课题以气候模式场增强为独立研究对象。

\subsubsection{可行性分析}

\paragraph{方法基础}
本项目从统计校正和等信息独立模型起步，逐步检验共享表征、地形条件与约束。前述先行试验已在公开资料上完成变量对齐、粗化算子构造、相同预算下的参数量匹配、两因素受控对照与聚合残差诊断，全流程在单张 A40 上约 28 分钟完成，表明计算规模与本课题资源相匹配，技术路线在原理上可行。既有时序建模经验可用于模型实现，公开降尺度研究可提供比较依据。所提出的科学问题可通过增量消融与独立验证回答，无需在一年内同时从头开发大型气候模式或全球基础模型。

\paragraph{资料与算力}
公开 CIESM 和再分析资料可用于先行研究，申请人已确认可稳定使用的 GPU 及软件环境能够支持区域分块训练。公开资料的变量频率和共同覆盖期仍需实际核查，因此资料获取与模型研究将分阶段推进，在样本条件明确前不预设性能结果。

\paragraph{组织安排}
申请人负责总体设计、模型与验证方法、进度管理和成果组织，重点承担目标场信息含量检验与 TerraSR-XJ 约束投影模块的实现。本课题组织新疆大学 7 名在读研究生参与研究，围绕资料处理、模型实验和集成评价三个方向分工协作。
\mbox{孔哲}、\mbox{张建宏}拟负责资料处理与质量控制，开展气象变量对齐、日历与累计量核验、粗化算子复算及台站隔离检查；\mbox{杨铄}、\mbox{孟星宇}、\mbox{桑乐}拟负责模型实验与消融，开展等信息独立基线、共享表征模型、地形因子及一致性约束的受控对照；\mbox{吕中原}、\mbox{于得武}拟负责集成评价与接口联调，开展台站独立验证、CIESM 文件处理、重建结果的粗化复核及运行记录整理。
各项任务由申请人统筹指导，研究生参与工作与科研培养相结合。本课题安排研究生劳务费 1.10 万元，按合计 22 人月、每人月 500 元测算，用于上述资料处理、模型实验与集成评价任务，并按实际参与工作量分配。

\subsubsection{研究风险与应对措施}

\paragraph{数据与频率风险}
若第一季度发现 CIESM 关键变量缺失或仅有月频资料，将及时寻求实验室提供合适版本，同时保留公开资料的可复现先行试验。月尺度结果只作为阶段性证据，不替代日尺度核心考核。针对目标场的插值天花板风险，本课题在建模前开展信息含量检验；未通过检验的变量退出重建任务，据实说明原因，且不进入性能汇总。若气温之外的三个变量在 CMFD 上同样未通过检验，则提出“气温重建 + 四变量一致性校正与偏差诊断”的调整方案，并与实验室确认其对研究目标和考核的影响。

\paragraph{技术风险}
若联合训练使某些变量持续退化，则采用等信息独立输出或降低共享深度；若软约束与极端结构冲突，则保留数学上成立的条件，减弱经验惩罚。调整依据来自固定验证集，修改后重新进行受控比较。若随机模型无法形成可校准的样本分布，则交付确定性核心模型及其误差说明，不对未经检验的概率结果作出承诺。

\paragraph{进度与应用风险}
实验室接口对接从第一季度开始，公开格式测试与接口联调分开推进。若正式接收延迟，将保留可复现兼容性结果，如实列示未完成的应用环节，并继续按约定联调。论文尽早完成主要设计和初稿，预留审稿修改时间；论文投稿、接收与发表的进度分别记录。

\paragraph{资源风险}
采用区域小块、轻量模型和少量种子重复控制训练规模，并记录 GPU 时间和存储占用。实验限定为新疆区域四类变量的日尺度重建，资源主要用于等信息对照、约束消融与集成验证。若资料规模或生成模型开销仍高于预期，则将概率扩展顺延至后续工作，核心目标不依赖尚未落实的额外算力。

% 迁移来源：NSFC_2026_CIESM-STSR/extraTex/3.2.工作条件.tex
\subsubsection{已具备的实验条件}

申请人已确认可稳定使用 NVIDIA A40 共 8 张和 NVIDIA GeForce RTX 3090 Ti 共 6 张。现有 Linux、Python、深度学习框架、Git 和 LaTeX 工具链能够支持数据处理、区域模型训练、版本管理与结果整理。实验将按实际调度分批开展，研究预算不以同时独占所有显卡为前提。

本课题主要使用公开 CIESM 历史模拟、ERA5、ERA5-Land 和公开地形资料，目前已具备通过公开渠道获取资料并开展方法验证的条件。具体文件、变量频率与覆盖年份将在项目开始后逐项核查；核心任务不依赖私有新能源场站数据或业务数值天气预报归档。

\subsubsection{尚缺少的实验条件及拟解决的途径}

目前尚需落实的仅有实验室合作人、内部 CIESM 数据及接口样例，相关对接事项与验收方式已在“利用研究基地的计划与落实情况”中列明。拟在第一季度对接数据格式、变量含义及实际应用需求，同时在公开资料上先完成处理流程。内部数据的访问与使用遵循实验室要求，在落实前仅作为拟争取条件；本课题四项研究内容均可基于公开资料完成，不以内部数据为前提。

公开渠道已能提供独立观测验证所需资料。中国气象数据网公开发布的地面气候资料日值产品可按其许可条款申请获取，本课题拟在新疆区域选取国家级台站，建立不参与训练的独立验证集；目标场之一的 CMFD 也由国家青藏高原科学数据中心公开发布\cite{He2020CMFD}。二者均无需本项目专属商业授权，台站独立验证因此纳入必做环节。由于 CMFD 生成使用了台站信息，验证集必须与参与 CMFD 生成的站点严格隔离，隔离清单将在资料核查阶段确定并公开。受资助规模所限，本课题缺少加密自动站、通量塔与商业气象数据，也未将其纳入预算。学院公共算力可申请作为补充资源，不预先计入固定配额。

\subsubsection{利用研究基地的计划与落实情况}

拟依托新疆大学现有科研环境开展算法研究，并与地球系统数值模拟教育部重点实验室围绕 CIESM 数据增强进行技术交流和应用对接。

\paragraph{使用的模式资料与配置}
拟使用 CIESM 参加 CMIP6 的 \texttt{historical} 试验输出，其大气分量水平分辨率约为 1$^{\circ}$ 量级，重建目标网格为 0.1$^{\circ}$，对应线性约 10 倍、面积约 100 倍的提升。第一季度资料核查将依据实际文件元数据，确定并记录模式版本、成员标识、网格名称、变量频率与文件校验值；先使用公开渠道可获取的资料，实验室内部版本按其要求申请。

\paragraph{对接事项与分工}
拟请实验室方指定一位合作人指导数据处理与接口对接，协助确认可用的 CIESM 版本与变量清单，提供接口样例与输出语义说明，并参与增强模块验收。本课题核心任务基于公开资料设计，实验室内部数据为拟争取的补充研究条件。

\paragraph{算法模块的提交与验收}
TerraSR-XJ 模块将以代码仓库、接口文档、运行样例与单元测试的形式提交实验室。验收时，在实验室指定的 CIESM 文件上完成读入、分块推断、细网格输出与粗化复核全流程，由合作人确认输出语义与运行记录。公开格式兼容性、实验室接口联调和实际应用的证据分别留存，仅在取得相应记录后报告相关工作的完成情况。

\paragraph{成果归属与署名}
本课题产出论文将标注地球系统数值模拟教育部重点实验室为署名单位，邀请实验室合作人共同署名，并在致谢中注明本课题受该实验室开放基金资助。区域试验数据产品将在对外发布前提交实验室审核，许可范围、归属与发布渠道按实验室要求确定；算法代码的开源范围也遵循实验室约定及各上游数据许可。若实验室对上述任一项另有规定，则以实验室规定为准。

% 迁移来源：NSFC_2026_CIESM-STSR/extraTex/3.3.承担项目情况.tex
\subsubsection{相关在研项目}

申请人正在承担且与本课题方法或资料处理有关的项目如下（截至 2026 年申报时点）。

\paragraph{融合多源数据大模型驱动的光伏发电功率预测方法研究}
自治区高校基本科研业务费科研项目—培育类，批准号 XJEDU2026P005，资助金额 5 万元，其中直接经费 3.9 万元、间接经费 1.1 万元，执行期 2026 年 1—12 月。申请人主持该项目，负责总体方案、模型研究和成果组织。该项目研究光伏功率，本课题研究 CIESM 气象场跨尺度重建，二者可共享时序建模经验；气候场增强工作或拟产出的数据不重复列为该项目已完成成果。

\paragraph{多模态融合技术研究及其在新疆产业集群中的应用}
自治区天池英才引进计划（青年博士人才），依据任务书无独立项目批准号，资助金额 50 万元，执行期 2025 年 2 月至 2028 年 2 月。申请人主持该项目，负责总体研究方案、关键模型设计和成果组织。相关多源信息融合经验可用于本课题的方法实施，本课题则以 CIESM 数据、物理范围与聚合一致性作为研究任务和评价对象。

\paragraph{数据驱动下基于互补型生成对抗网络与神经统计模型的公路沙害智能预警研究}
国家自然科学基金地区项目，批准号 62466056，总资助金额 32 万元，执行期 2025 年 1 月至 2028 年 12 月，基金委系统登记单位为新疆大学。申请人作为参与者，负责数据收集整理，相关经验可用于本课题的区域资料组织。本课题研究气候模式数据增强，不重复研发该项目的公路沙害预警模型。

上述在研项目积累的算力与存储条件可用于本课题的数据处理和模型训练，本课题不另行购置设备。本次申请的 \CiesmValue{total-budget} 万元用于 CIESM 对接差旅、专题交流、成果发表与传播、专家咨询及研究生劳务，其中劳务费 1.10 万元用于研究生参与本课题相关任务。三项在研项目分别研究光伏功率、产业集群多模态融合与公路沙害预警，均不以气候模式场跨尺度重建、聚合一致性约束或模式数据集增强为任务，本课题成果不重复计入上述任一项目的考核。
